\chapter{The IR}
\label{ch:ir}

The \term{IR} (intermediate representation) is a YAML file that holds the whole design
after BAG has resolved it. Every generated file comes from it, so the IR is the contract
between the workbook and the output. \cmd{bag ir} writes it:

\begin{center}
\cmd{bag ir spec.xlsx -o spec.ir.yaml}
\end{center}

\section{What the IR holds}

The IR holds what you typed and what BAG derived from it:

\begin{itemize}
\item the project, the clocks and resets, and the wrapper settings, as typed;
\item every fabric with its ports in port order, its ID widths, its address rules (the
  decode table, with \sheet{Access_Rules} and aliases applied) and its decode settings per
  master;
\item every master, slave and bridge, with the fabric it is on;
\item every adapter BAG inserts (protocol, width, clock domain, ID width, FIFO, register
  slice), with its instance name and the sub-wrapper it is placed in;
\item the memory map, the port options, the user ports and the hierarchy;
\item a \code{provenance} section: the version of BAG, the workbook it came from and its
  SHA-256 hash.
\end{itemize}

This is the IR of the small design of chapter~\ref{ch:first}: one APB4 fabric
\code{apb} with the master \code{m0} and the slave \code{s0}.

\lstinputlisting[style=yaml]{examples/files/apb4_min.ir.yaml}

Keys and lists come in a fixed order, and numbers in a fixed form (addresses as
\code{'0x0000_0800'}), so the same design always gives the same file.

\section{The IR as input}

\cmd{bag generate} and \cmd{bag flow} take an IR as well as a workbook. (\cmd{bag check}
and \cmd{bag ir} read workbooks only.) An IR may have been edited by hand, so BAG checks
it as strictly as a workbook: it converts the IR back to a workbook and runs every check
on that. The summary line says so. A finding names a sheet and row of that converted
workbook, which is not written to disk; use the names and values in the message to find
the place in the IR.

\lstinputlisting[style=console]{examples/console/ir-input.txt}

In \file{apb4_min_edited.ir.yaml} the size of the region \code{regs} of \code{s0} is
\code{0x0000_0C00}, which is not a power of 2. The second error follows from the first:
the alias of \code{m0} points into that region.

\section{The schema version}

The IR starts with \code{schema:}, the version of the workbook and IR schema. This
version of BAG reads schema 0.6, and the Project sheet of the workbook names the same
version. BAG refuses an IR of another version with a message such as
\code{x.ir.yaml is schema 0.5; this bag reads schema 0.6. Update its workbook
(docs/schema-changelog.md), then write the IR again with bag ir}, and an IR with a key it
does not know with \code{x.ir.yaml is not a valid IR: project.colour: Extra inputs are
not permitted}. The file \file{docs/schema-changelog.md} of the repository lists what
changed in each version.

\section{Workbook or IR?}

Keep the workbook as your source. It is what a person reads and edits, and the checks
report against its rows. Keep the IR under version control next to it: a review of an
IR diff shows exactly what a workbook change did to the design, including the derived
parts such as inserted adapters and decode tables.
